\appendix

\section{Appendix: Implementation Details}
During sampling, we perform DDIM sampling \citep{ddim} with $50$ denoising steps, setting DDIM eta to 0. 
%
The inference resolution is fixed at $256\times256$ pixels. The scale of the classifier-free guidance is set to $15$.

For quantitative comparison, we generate a total of $2048$ videos for each longer inference method, utilizing $512$ prompts (from a standard evaluation paper EvalCrafter~\citep{liu2023evalcrafter}) and initializing with $4$ random initial noises. For calculating the FVD and KVD, we use videos generated by direct inference with training length ($16$ frames) as the reference set. To align the video length, we cut each long video ($64$ frames) into $4$ segments with $16$ frames. Correspondingly, we generate the same number ($8192$) of short videos with direct inference.

In the user study, we mixed our generated videos with those generated by the other three baselines. A total of $27$ users were asked to pick the best one according to the content consistency video quality and video-text alignment, respectively.

\begin{figure*} %[t]
%
\centering
\includegraphics[width=0.98\linewidth]{Figures/fig_motion.pdf}
\caption{FreeNoise can produce three types of videos exhibiting significant movement: (a) the lens moving with the subject, (b) the subject moving off the screen, and (c) the subject moving within the screen.
}
\label{fig:motion}
\end{figure*}

\section{Appendix: Case Analysis of Significant Movement}
Videos with significant movement can be mainly divided into three types: 
%

\noindent\textbf{Lens Moving with the Subject.}
For videos of this type, the position of the subject does not change much and the movement is shown through the regression of the background (Figure~\ref{fig:motion}(a)).

\noindent\textbf{Subject Moving off the Screen.}
For videos of this type, the subject will move off the screen (Figure~\ref{fig:motion}(b)). However, the subject will suddenly appear again due to semantic constraints.

\noindent\textbf{Subject Moving within the Screen.}
For videos of this type, the subject will move within the screen. Due to the size limitation of the screen, the subject will turn around (Figure~\ref{fig:motion}(c)). However, the current pretrained model behaves unnaturally when turning (even for original inference).

As shown in Figure~\ref{fig:motion}, FreeNoise is able to produce all three types. These three types are automatically determined during the inference stage based on the sampled random noises and the given prompt.
%
% All three types of generated videos are available at the anonymous website: \href{https://free-noise.github.io}{https://free-noise.github.io}.

\section{Appendix: Limitation Discussion}

As repeated locally shuffled noises are used, FreeNoise has a weakening effect on introducing new content to the video as the length increases. In some cases, the displacement of the subject is limited due to this weakening effect of FreeNoise. However, FreeNoise does not obliterate motion variation or thoroughly fix the spatial structure of objects, such as an object moving from left to right on the screen.

Currently the base model (inference without FreeNoise) only works well with videos of the lens moving with the subject and still struggles to deal with the videos of the subject moving off/within the screen effectively. Therefore, the performance of FreeNoise is also constrained by the base model. We look forward to applying FreeNoise to more powerful video models in the future.
